\begin{frame}{손으로 한 번: Peer-Review rubric을 직접 써보기}
  \small
  \begin{itemize}
    \item 관점을 한 단계 더 뒤집어 \textbf{``너가 채점자라면 무엇을
      봐야 하는지'' rubric을 직접 쓰는 연습}.
    \item 5개 항목을 그대로 두고, 각 항목에 대해 \textbf{두 칸}을 채운다:
      (1) \textbf{나쁠 때는 어떤 모습인가}(구체적 나쁜 패턴), (2)
      \textbf{그걸 시험하는 질문은 무엇인가}(16.2의 반박 가능한 질문).
  \end{itemize}
  \vskip0.4em
  \begin{tabular}{@{}p{0.52\textwidth}l@{}}
    \toprule
    \textbf{나쁠 때 (구체적 패턴)} & \textbf{시험하는 질문 (반박 가능)} \\
    \midrule
    보고서에 \textbf{최고 점수 시드 하나만} 실려 있음. 학습 곡선이
    매끄럽고 불안정 구간 설명이 없는데, 학습 로그(코드 출력)를 보면
    300 iteration 중 \textbf{150$\sim$200 구간에서 리턴이 갑자기
    내려간} 기록이 있음 & ``로그의 150$\sim$200 iteration 리턴 추락을
    보고서의 학습 곡선에서 \textbf{왜 볼 수 없나요?} 그 구간을 자르면
    곡선이 매끄럽혀지나요?'' — \textbf{실제로 다시 돌려봐야} 답할 수
    있는 질문 \\
    \bottomrule
  \end{tabular}
  \vskip0.6em
  {\footnotesize 이 연습의 채점 기준 = 16.2 ``구체적인가'' 3단계
  그대로: ``나쁨 = 불성실''은 \textbf{1점}(감상평), 그날 발표실에서
  \emph{실제로 낚일 수 있는 패턴}을 쓴 것은 \textbf{3점}.
  ``나쁠 때의 모습''을 쓸 수 있으려면 ``좋은 답이란''을 이미 정확히
  이해해야 하므로, 이 30분 연습(나머지 4행 = 팀 회의 과제)은
  \textbf{학기 전체를 가장 빨리 재검증하는 방법}.}
\end{frame}
